\cotefile{156-9}
\scanpage{134}
\mspage{5}
% archive sheet number: 133

\uline{Définition} (provisoire) \unsure{soient|soit} $A$\unsure{,|.} $A'$ deux ateliers. Un \deleted{\unsure{un}} homomorphisme (logique, sans \unsure{connotation} \unsure{géométrique|géometrique} fixée, pour le moment) de $A$ dans $A'$, est une application \deleted{de $A$}
\[
f:A\longrightarrow\inserted{\widehat{A}':=}A'\amalg\{\varnothing\}\quad(\hookrightarrow\unsure{F|\mathcal{F}}(A'))
\]
satisfaisant les trois conditions, pour $X,Y\in A$
\[
\begin{aligned}
\text{\unsure{i)|a)}}\quad &X\trianglelefteq Y\Longrightarrow f(X)\trianglelefteq f(Y)\\
\text{\unsure{ii)|b)}}\quad &X\overset{\circ}{\ll}Y\Longrightarrow f(X)\overset{\circ}{\ll}f(Y)\\
\text{\unsure{iii)|c)}}\quad &X\mathrel{|\circ|}Y\Longrightarrow f(X)\mathrel{|\circ|}f(Y).
\end{aligned}
\]
% margin note written sideways
% slanted
\margin{Mais la condition \unsure{d)} de \unsure{GF|CF} $\overline{\text{\uline{VIII}}}$ \unsure{p.~27|n'est|n'en} \unsure{n'est|vaut|tient|valable} \unsure{pas|pour} \unsure{compte|comptée} pour \unsure{le} moment ! \unsure{!)|i)}}

Pour donner un sens à cette définition, faut donner un sens aux relations
\[
X'\trianglelefteq Y',\quad X'\overset{\circ}{\ll}Y',\quad X'\mathrel{|\circ|}Y'\qquad\text{pour }X',Y'\in A'\amalg\{\varnothing\},
\]
\unsure{des symboles}, quand l'un ou l'autre des deux termes est \deleted{$\{\varnothing\}$} égal à $\varnothing$. Les bonnes conventions sont pour $\forall\,X'\in A'$
\[
\text{a)}\quad X'\trianglelefteq\varnothing,\quad
\deleted{\unsure{\forall X'}}\quad
(\varnothing\trianglelefteq X'\ \deleted{\unsure{\Longleftrightarrow}}\ \text{tjs vrai}
\]
\margin{(\uline{NB} $\trianglelefteq$ sur $\widehat{A}'\deleted{\unsure{\amalg\{\varnothing\}\subset A}}$ n'est \unsure{pas} une relation d'ordre !)}
\[
\text{b)}\quad\varnothing\overset{\circ}{\ll}X',\qquad
(X'\overset{\circ}{\ll}\varnothing\Longleftrightarrow X'=\varnothing)
\]
\margin{\uline{NB} $\overset{\circ}{\ll}$ sur $\widehat{A}'$ est une rel. d'ordre admettant $\varnothing$ comme plus petit él.}

\scanpage{135}
\mspage{6}
% archive sheet number: 134

\[
\text{c)}\quad X'\mathrel{|\circ|}\varnothing\qquad\text{tjs satisfaite.}
\]

L'introduction de $\widehat{A}'$, et les conventions précédentes, sont inspirées par l'exemple $2^\circ)$ précédent.

\uline{Noyau} de $f$ :
\[
\operatorname{Ker} f\overset{\mathrm{dfn}}{=}\{X\in A\mid f(X)=\varnothing\}
\]
Alors b) implique que $\operatorname{Ker} f\subset A$ est fermé pour $\overset{\circ}{\ll}$. Du coup on s'attend que ce soit dans $\Sigma_A$, mais on doit sans doute pour cela imposer des conditions \unsure{supplémentaires} sur $f$ (en tous cas c'est OK dans l'exemple $2^\circ$ et ses variantes).

Partons d'une partie $\unsure{N}\subset A$, fermée pour $\overset{\circ}{\ll}$, est-ce le noyau d'un hom. d'atelier, et ce dernier peut-il se choisir canoniquement ? Plus précisément, peut-on trouver un atelier $A/N$ (opération d'\uline{excision})

\scanpage{136}
\mspage{7}
% archive sheet number: 135

et un hom
\diagram{\begin{tikzcd}[ampersand replacement=\&,column sep=small]
A\&\overset{p}{\unsure{{-}|\longrightarrow}}\& A/N
\end{tikzcd}}
ayant le noyau $N$ (ou $\supset N$ ?), qui soit \uline{universel}, en ce sens que tout hom d'ateliers
\diagram{\begin{tikzcd}[ampersand replacement=\&,column sep=small]
A\&\overset{f}{\unsure{{-}|\longrightarrow}}\& A'
\end{tikzcd}}
ayant un noyau $\supset N$, se factorise par $p$ de façon unique. (\uline{NB} La notion de \uline{composition} est claire, en interprétant les hom. d'ateliers comme des applications ensemblistes
\diagram{\begin{tikzcd}[ampersand replacement=\&,column sep=small]
A\amalg\{\varnothing\}\&\overset{\widehat{f}}{\unsure{{-}|\longrightarrow}}\& A'\amalg\{\varnothing\}
\end{tikzcd}}
transformant $\{\varnothing\}$ en $\{\varnothing\}$, et \deleted{\unsure{un hom}} \deleted{\unsure{que}} dont la restriction à $A$ a les propriétés préc. \unsure{cela revient} à dire que $\widehat{f}$ est compatible à $\trianglelefteq,\overset{\circ}{\ll},\mathrel{|\circ|}$ sur $\widehat{A}$ !
% margin note written sideways
% slanted
\margin{\unsure{commencer} par ça !}

L'exemple type est celui où
\[
A=\unsure{\operatorname{Fig\acute{e}lem}|\operatorname{Fig\acute{e}l\acute{e}m}}(E),\quad N=\{X\in A\mid |X|\subset E\setminus E')
\]
\[
A/N=\unsure{\operatorname{Fig\,elem}}(E'),\quad\text{mais}
\]
\unsure{on veut} \unsure{trouver|toujours} \inserted{dans ce cas} \unsure{interprété|interpréter} comme quotient de $A\setminus N$

\scanpage{137}
\mspage{8}
% archive sheet number: 136

À supposer qu'une telle construction existe dans un cas tant soit peu général, elle s'annonce un \unsure{tantinet} délicate. Pour fixer les idées, il serait bon \deleted{\unsure{tjs|de}} au moins de regarder déjà les cas suivants

$1^\circ)$ Cas d'un atelier ensembliste

\deleted{\unsure{$2^\circ)$}} Déjà \unsure{ça} n'a pas l'air évident, surtout pas quand on ne suppose pas que l'atelier soit modéré, et $\unsure{\eqsim|\simeq|=}$ pour des $N$ très particuliers comme p.ex. $\operatorname{Omb}^{\unsure{\bullet|\circ}}F$, $F$ une figure, ou $\operatorname{Supp}F$.

$2^\circ)$ Cas d'une maquette $A$, $N$ partie quelconque de $A$ (NB $x\overset{\circ}{\ll}y\Leftrightarrow x=y$ !). Là, on doit trouver la maquette $A\setminus N$ simplement. Ce n'est \unsure{pas} \unsure{sorcier} cette fois !

\par\noindent\hspace*{0.3\linewidth}\rule{0.25\linewidth}{0.4pt}\par

Pour la relation avec le \unsure{formalisme} des supports, s'inspirer éventuellement de \unsure{GF|CF} $\overline{\text{\uline{VIII}}}$ p.~70 ff. Là je n'ai pas la pêche pour, et laisse tomber pour le moment.

\scanpage{138}
\mspage{9}
% archive sheet number: 137

\uline{Subdivisions}

Soient $A$ un \deleted{\unsure{organisme}} \inserted{atelier}, $\unsure{\mathcal{F}|\mathscr{F}}$ l'ensemble des $A$-figures, et
\[
G\ll F
\]
dans $\mathcal{F}$. On a donc une application \uline{croissante}
\diagram{\begin{tikzcd}[ampersand replacement=\&,row sep=small]
\unsure{\varphi|\phi}:\&\widetilde{G}\arrow[r]\arrow[d,equals]\&\widetilde{F}\arrow[d,equals]\\
\&G\&F
\end{tikzcd}}
telle que pour $Y\in\widetilde{G}$, $X=\varphi(Y)$, on ait
\[
Y\overset{\circ}{\ll}X
\]
Pour dire que $G$ est une \deleted{\unsure{raffinement}} \uline{subdivision} de $F$, nous exigerons \deleted{d'abord} que

(i) $\forall\,X\in\widetilde{F}$, \deleted{posant} considérant $\varphi^{-1}(X)\subset\widetilde{G}$,
$\varphi^{-1}(X)\overset{\circ}{\ll}X$, et $\operatorname{supp}^{\circ}\varphi^{-1}(X)\subset X^{\circ}$, on a
\deleted{on a}\inserted{donc}
\[
\deleted{\operatorname{supp}^{\circ}\varphi^{-1}(X)}
\]
\[
\operatorname{supp}^{\circ}\varphi^{-1}(X)=X^{\circ}
\]
% margin note written sideways
% slanted
\margin{Il valait mieux ici partir \unsure{d'une|d'un} \deleted{\unsure{subdivision|atelier}} \inserted{\unsure{figure}} \unsure{élém.} $F$ \unsure{donné|donnée} \unsure{et sa topologie}.}

En \unsure{admettant|admet.} \uline{\unsure{Atmag|Atomag} supp$^2$}, on voit que ceci équivaut à chacune des deux conditions

(ii)
\diagram{$\begin{array}{ccc}
\operatorname{supp}G&=&\operatorname{supp}F\\
\Vert&&\Vert\\
(\operatorname{supp}^{\unsure{\bullet|\circ}}G)&&(\operatorname{supp}^{\circ}F)
\end{array}$}
ou encore

(iii) Il n'existe pas $Z\in A$ \unsure{avec} $Z\ll F$, et $Z\mathrel{|\unsure{\circ|\bullet}|}\deleted{\unsure{F}}G$.

\scanpage{139}
\mspage{10}
% archive sheet number: 138

Quand $A$ est divisible, \deleted{\unsure{on}} $Z\in\unsure{L|\mathcal{L}}$, et on trouve la forme

\deleted{\unsure{(ii'')|(iii')}}\inserted{(iv)} Il n'existe pas $\unsure{x|\alpha}\in L$, avec $x\in\operatorname{Omb}(F)$ (i.e. $x\ll F$), et $x\mathrel{|\circ|}G$.

En fait, les implications
\[
\text{(i)}\Longrightarrow\text{(ii)}\Longrightarrow
\deleted{\text{\unsure{(ii')}}}\text{(iii)}\Longrightarrow
\deleted{\text{\unsure{(ii'')}}}\text{(iv)}
\]
sont \deleted{\unsure{est}} triviales, et $(\deleted{\unsure{\mathrm{iii}}}\mathrm{iv})\Longrightarrow\text{(i)}$ (moyennant divisibilité) \unsure{resp.} $\text{(iii)}\Longrightarrow\text{(i)}$, grâce à \unsure{Atom.g|Atmag|Atomag} supp$^2$.

Signalons enfin une autre condition équivalente, moyennant \uline{divisibilité} :

\deleted{\unsure{(ii quater)|(ii''')}}\inserted{(v)} $G$ est maximal parmi les raffinements de $F$, ordonnés par $\deleted{\unsure{\ll}}\subset$.
% margin note written sideways
% slanted
\margin{\uline{NB} Ces conditions impliquent déjà que $\unsure{\varphi|\phi}$ est \unsure{surjective}.}

\deleted{Mais}\inserted{Quand} ces conditions \deleted{\unsure{n'impliquent pas}}\inserted{sont satisfaites}, on dira \deleted{\unsure{pas encore}} que $G$ est une \uline{subdivision de} $F$. \deleted{(Comme précédemment)}

% The following entire block is crossed through with diagonal strokes.
\deleted{mais je voudrais la condition}

\deleted{(ii) La top. de $\widetilde{F}$ (associée à sa structure d'ordre) est \uline{quotient} de celle de $\widetilde{G}$, ou encore, si $X,X'\in\widetilde{F}$, pour qu'on ait $X\leq X'$ il faut (et suffit) qu'il existe $Y\in\varphi^{-1}(X)$, $Y'\in\varphi^{-1}(X')$, avec $Y\leq Y'$.}
% margin note written sideways
% slanted
\margin{\unsure{est-elle indépendante} de l'exigence \unsure{?} \uline{Voire !}}

\deleted{Cela signifie aussi, dans le cas $A=\unsure{\operatorname{Fig\acute{e}lem}|\operatorname{Fig\acute{e}l\acute{e}m}}(\unsure{E|L}$}

\scanpage{140}
\mspage{11}
% archive sheet number: 139
% The entire text block is crossed through with two diagonal strokes.

\deleted{disons, que $F$ se déduit de $G$ par ``regroupement de termes'' ${\unsure{Y|\dot{Y}}}_j^{\circ}$ (pour former des $X_i^{\circ}$) sans plus -- et que la relation d'ordre sur $I$, déterminant les ``\unsure{strates} \unsure{fermees|fermées}'', se déduit simplement de celle de $J$ par passage au quotient.}
